% part: intuitionistic-logic
% chapter: propositions-as-types
% section: types

\documentclass[../../../include/open-logic-section]{subfiles}

\begin{document}

\olfileid{int}{pty}{typ}

\olsection{টাইপ}

$(MN)$-এর মতো প্রমাণপদ একাই তার
সংশ্লিষ্ট প্রমাণটিকে অনন্যভাবে নির্ধারণ
করে না। কারণ মুক্ত চলরাশিগুলি, অর্থাৎ
স্বতঃসিদ্ধের প্রমাণপদগুলি, সংশ্লিষ্ট
স্বতঃসিদ্ধের !!{formula}s নির্দিষ্ট করে
না। কিন্তু প্রসঙ্গ~$\Gamma$-তে সেগুলি
নির্দিষ্ট করলে প্রমাণটি \emph{অনন্যভাবে
নির্ধারিত হয়}।

\begin{defn}
প্রমাণপদ~$N$-এর \emph{প্রসঙ্গ} হলো এমন
সমষ্টি~$\Gamma$, যাতে $N$-এর প্রত্যেক
মুক্ত চলরাশির জন্য $\Gamma$-তে ঠিক
একটি $x : !A$ যুগল আছে।
\end{defn}

সংজ্ঞায় বলা হয়নি যে $\Gamma$-তে
$x:!A$ যুগল \emph{কেবল} $N$-এর মুক্ত
চলরাশির জন্যই থাকতে হবে। যেমন,
$\Gamma = \{x:!A, y:!B, z:!C\}$
হলো $\pair{x}{y}$-এর একটি প্রসঙ্গ।

\begin{defn}\ollabel{defn:type}
  ধরি, $\Gamma$ হলো~$N$-এর প্রসঙ্গ।
  $\Gamma$-র সাপেক্ষে $N$-এর \emph{টাইপ}
  আরোহীভাবে নিম্নমতো সংজ্ঞায়িত:
  \begin{enumerate}
  \item $x$-এর টাইপ $!A$ তখনই এবং কেবল
  তখনই, যখন $x:!A \in \Gamma$।
  \item $x:!A, \Gamma$-র সাপেক্ষে $N$-এর
  টাইপ~$!B$ হলে $\lambd[\typeof{x}{!A}][N]$-এর
  টাইপ $!A \lif !B$।
  \item $N$-এর টাইপ $!A \lif !B$ এবং $M$-এর
  টাইপ~$!A$ হলে $(NM)$-এর টাইপ~$!B$।
  \item $N$-এর টাইপ $!A$ এবং $M$-এর
  টাইপ~$!B$ হলে $\pair{N}{M}$-এর টাইপ
  $!A \land !B$।
  \item $N$-এর টাইপ~$!A_1 \land !A_2$
  হলে $\proj{i}{N}$-এর টাইপ~$!A_i$।
  \item $N$-এর টাইপ $!A$ হলে
    $\inj[!B]{i}{N}$-এর টাইপ $i=1$-এ
    $!A \lor !B$, আর $i=2$-এ
    $!B \lor !A$।
  \item $M$-এর টাইপ $!A \lor !B$ হলে,
  $x:!A, \Gamma$-র সাপেক্ষে $N_1$-এর
  টাইপ $!C$ এবং $y:!B,\Gamma$-র
  সাপেক্ষে $N_2$-এর টাইপ~$!C$ হলে,
  $\dcase{M}{x}{N_1}{y}{N_2}$-এর টাইপ~$!C$।
  \item $N$-এর টাইপ~$\lfalse$ হলে
    $\abort{!A}{N}$-এর টাইপ~$!A$।
  \end{enumerate}
\end{defn}
% BN-SRC-621 TARGET-CORRECTION-BEGIN
\par\smallskip\noindent\textnormal{\small\textbf{উৎসসংশোধন (BN-SRC-621):} যোগ-অন্তর্ভুক্তি ধারায় উৎসের \(\inj[!B]{i}{!M}\)-এ অনির্ধারিত \(M\) এবং অপ্রয়োজনীয় বিস্ময়চিহ্ন ছিল; এই ধারায় বাঁধা প্রমাণপদ \(N\), তাই \(\inj[!B]{i}{N}\) করা হয়েছে।}
% BN-SRC-621 TARGET-CORRECTION-END
% BN-SRC-622 TARGET-CORRECTION-BEGIN
\par\smallskip\noindent\textnormal{\small\textbf{উৎসসংশোধন (BN-SRC-622):} যোগ-নিষ্কাশনে পূর্বে বাঁধা শাখাচলরাশি \(x,y\), অথচ উৎসের \(\dcase\)-এ \(x_1,x_2\) ছিল; সমান্তরাল tN3i বিধির মতো \(\dcase{M}{x}{N_1}{y}{N_2}\) করা হয়েছে।}
% BN-SRC-622 TARGET-CORRECTION-END

\begin{prop}
প্রমাণপদ~$N$-এর যেকোনো প্রসঙ্গ~$\Gamma$-র
সাপেক্ষে $N$-এর ঠিক একটি টাইপ আছে।
\end{prop}

\begin{proof}
  $N$-এর গঠনের উপর আরোহে।
\end{proof}

$\Gamma$-র সাপেক্ষে $N$-এর টাইপ~$!A$
হলে লিখি $\Gamma \Sequent N : !A$।
\olref{defn:type}-এর ধারাগুলি পরিচিত:
একটি ছোট পার্থক্য ছাড়া এগুলি ঠিক
\Log{tN2i}-এর বিধি—এখানে প্রসঙ্গ~$\Gamma$
সর্বত্র একই। সংজ্ঞাটিকে \Log{tN3i}
কলনে রূপ দেওয়া যায়, যার বিধি
\olref{tab:tN3ip}-তে আছে।
% BN-SRC-623 TARGET-CORRECTION-BEGIN
\par\smallskip\noindent\textnormal{\small\textbf{উৎসসংশোধন (BN-SRC-623):} এখানে আমদানিকৃত সারণি tN3i-র; উৎসের \texttt{tab:tN2ip} উল্লেখটি আগের tN2i সারণিতে যায়। পূর্বে স্থির স্বতন্ত্র \texttt{tab:tN3ip} লেবেলের সঙ্গে মিলিয়ে উল্লেখ সংশোধিত।}
% BN-SRC-623 TARGET-CORRECTION-END

\subfile{rules-tN3}

প্রমাণপদগুলি $\lambd$-ক্যালকুলাসের একটি
প্রকরণ, \emph{গুণফল, যোগফল ও শূন্য টাইপযুক্ত
$\lambd$-ক্যালকুলাস}-এর পদ। চিরায়ত
$\lambd$-ক্যালকুলাসে কেবল চলরাশি,
\emph{$\lambd$-বিমূর্তন}~$\lambd[x][N]$
এবং প্রয়োগ~$(MN)$ থেকে পদ গঠিত হয়;
টাইপযুক্ত $\lambd$-বিমূর্তন
$\lambd[\typeof{x}{!A}][N]$-এর
$!A$-টীকা সেখানে নেই। $\lambd$-ক্যালকুলাস
গণনার একটি টিউরিং-সম্পূর্ণ মডেল,
অর্থাৎ প্রত্যেক আংশিক গণনাযোগ্য অপেক্ষক
এতে একটি পদ দিয়ে উপস্থাপন করা যায়।
প্রমাণপদ ও টাইপ-বিধিতে নির্ধারিত
$\lambd$-ক্যালকুলাসের প্রকরণটি গণনার
সীমাবদ্ধ মডেল, তবে মূল ধারণাগুলি একই।
টাইপ আরোপেই সীমা: কেবল
\emph{সঠিকভাবে টাইপযুক্ত} পদ, অর্থাৎ
প্রসঙ্গ~$\Gamma$-র সাপেক্ষে টাইপ আছে
এমন প্রমাণপদই বৈধ। যেমন,
$\lambd[x][(xx)]$ চিরায়ত টাইপবিহীন
$\lambd$-ক্যালকুলাসে বৈধ, কিন্তু
$\lambd[\typeof{x}{!A}][(xx)]$
কোনো~$!A$-এর জন্যই সঠিকভাবে
টাইপযুক্ত নয়। কারণ $(xx)$-এর টাইপ
থাকতে হলে প্রথম~$x$-এর টাইপ
$!A \lif !B$ এবং দ্বিতীয়~$x$-এর
টাইপ~$!A$ হতে হবে; অথচ $!A$ হলো
$!A \lif !B$-এর প্রকৃত উপ-!!{formula},
তাই তা অসম্ভব।

স্বজ্ঞাতভাবে, কোনো টাইপের পদ যে ধরনের
বস্তুর নাম দেয়, সেই ধরনই তার টাইপ।
তাই $!A \land !B$ টাইপের পদ এমন
একটি যুগল বেছে নেয় যার প্রথম উপাদানের
টাইপ~$!A$, দ্বিতীয়টির টাইপ~$!B$;
অর্থাৎ সেটি গুণফল~$!A \times !B$-এর
!!a{element}। $!A \lif !B$ টাইপের
পদ $!A$ টাইপের বস্তু থেকে $!B$
টাইপের বস্তুতে একটি অপেক্ষক।
$!A \lor !B$ টাইপের পদ $!A$
টাইপের অথবা $!B$ টাইপের কিছু,
সঙ্গে কোনটি তা-ও জানায়। এটি
$!A$ ও $!B$ সমষ্টির বিযুক্ত যোগফলের
মতো, যাকে লেখা যায় $\{\tuple{i,x}
: i = 1 \text{ বা } 2, x \in !A \text{ যদি } i = 1, x \in !B \text{ যদি
} i=2\}$। $\lfalse$ টাইপটি শূন্য
সমষ্টি: কোনো পদের এই টাইপ থাকলে
তা কোনো বস্তুর নাম দিতে পারে না।
স্বাভাবিক নিষ্পাদন, \Log{tN3i}
সহ, সঙ্গতিপূর্ণ বলে খোলা অনুমান
ছাড়া $\lfalse$-এর কোনো !!{proof}
নেই। তাই মুক্ত চলরাশিহীন কোনো
বদ্ধ প্রমাণপদের টাইপ~$\lfalse$ নয়।
% BN-NORM-180 TARGET-NORMALIZATION-BEGIN
\par\smallskip\noindent\textnormal{\small\textbf{ভাষাগত স্বাভাবিকীকরণ (BN-NORM-180):} বিযুক্ত যোগফলের সমষ্টি-সংজ্ঞায় উৎসের ইংরেজি \texttt{or/if/if} যথাক্রমে \texttt{বা/যদি/যদি} করা হয়েছে; সূচক, সদস্যতা, টাইপ ও শর্তের সত্যার্থ অপরিবর্তিত।}
% BN-NORM-180 TARGET-NORMALIZATION-END

টাইপযুক্ত $\lambd$-ক্যালকুলাসের
ক্রিয়কগুলি যে পদে প্রয়োগ করা হয়
সেগুলি নির্দিষ্ট ধরনের বস্তুর নাম
দিলে, ফলিত পদও স্বজ্ঞাতভাবে
যথোচিত ধরনের বস্তুর নাম দেয়।
যেমন, “$N_1 : !A$ এবং $N_2 : !B$
হলে $\pair{N_1}{N_2} : !A \land !B$”
বলছে: $N_1$ টাইপ~$!A$-এর এবং
$N_2$ টাইপ~$!B$-এর কোনো বস্তুর
নাম দিলে, $\pair{N_1}{N_2}$
টাইপ $!A \land !B$-এর কোনো
বস্তুর নাম দেয়।

\end{document}
